\scanpage{12}
\textbf{448 — Série A}\hfill\textbf{C. R. Acad. Sc. Paris, t. 279 (16 septembre 1974)}

\rule{\linewidth}{0.4pt}

\textsc{Proposition 3}. — \emph{Tout sous-Module et tout Module quotient d'un Module Semi-Simple
est encore Semi-Simple.}

\emph{Preuve}. — Elle est analogue à celle rencontrée dans le cas des modules semi-simples
non gradués.

\textsc{Définition 3}. — On appelle Radical de A et on note $\operatorname{Rad}(\mathrm A)$ l'intersection des Idéaux
à gauche maximaux de A. Le fait que $\operatorname{Rad}(\mathrm A)$ soit égal à l'intersection des annulateurs
des A-Modules Simples montre que c'est un Idéal bilatère. (Un élément $\alpha$ de A appartient
à l'annulateur de tout A-Module Simple $\Leftrightarrow$ $\alpha$ est l'annulateur de tout élément homogène
d'un A-Module Simple $\Leftrightarrow$ $\alpha$ appartient à tous les Idéaux à gauche maximaux de A.)

\textsc{Proposition 4}. — \emph{Soit $k$ un corps commutatif. Une $k$-algèbre Simple de dimension finie
est Semi-Simple.}

\emph{Preuve}. — Soit $\mathbf r$ un élément minimal de l'ensemble des intersections finies d'Idéaux à
gauche minimaux de A. Si $\mathbf m$ est un Idéal à gauche maximal de A, on a $\mathbf m\cap\mathbf r=\mathbf r$ d'où
$\mathbf r\subseteq\mathbf m\Rightarrow\mathbf r\subseteq\operatorname{Rad}(\mathrm A)$ ce qui entraîne que $\mathbf r=0$ vu que $\operatorname{Rad}(\mathrm A)$ est un idéal bilatère. Il
existe donc une famille finie $(\mathbf m_i)_{1\le i\le n}$ d'Idéaux à gauche maximaux dont l'intersection
est nulle. L'anneau A s'injecte dans le produit des $\mathrm A/\mathbf m_i$ et $\unsure{\mathrm A\,s|\mathrm A_s}$ est par conséquent semi-simple.

\textsc{Proposition 5}. — \emph{Une $k$-Algèbre Simple} A \emph{de dimension finie sur un corps $k$ est telle
que tout A-Module soit isotypique de type} A.

\emph{Preuve}. — D'après la proposition précédente, nous savons déjà que A est semi-simple.
Les Idéaux bilatères de A sont des sous-modules de $\unsure{\mathrm A\,s|\mathrm A_s}$ stables par multiplication à droite,
donc stables par le commutant de $\unsure{\mathrm A\,s|\mathrm A_s}$ qui n'est autre que l'anneau des homothéties à
droite de A ($\mathrm C\in\mathcal L_{\mathrm A}(\unsure{\mathrm A\,s|\mathrm A_s})$, $\mathrm C(y)=\mathrm C(\unsure{y_1|y1})=y\,\mathrm C(1)$). Ce sont donc des sommes de composants
isotypiques de $\unsure{\mathrm A\,s|\mathrm A_s}$ [(${}^1$), p. 34]. En particulier, les composants isotypiques de $\unsure{\mathrm A\,s|\mathrm A_s}$
sont les Idéaux bilatères minimals de A et la proposition en résulte.

2. \textsc{Théorème de Skolem-Nother}.

\textsc{Théorème 1}. — \emph{Soit $k$ un corps commutatif,} A \emph{une $k$-Algèbre Centrale Simple,} B \emph{une
$k$-Algèbre Simple, $f$ et $g$ deux homomorphismes de} B \emph{dans} A. \emph{Alors il existe un élément
$s\in\mathrm U(\mathrm A_0)$ (groupe des éléments inversibles de $\mathrm A_0$) qui est tel que pour tout $b\in\mathrm B$, on ait}
\[
f(b)=s^{-1}g(b)s.
\]

\emph{Preuve}. — On suppose d'abord que $\mathrm A=\operatorname{END}(\mathrm V)$ où V est un A-Module. On munit V
de deux structures de B-Module :
\[
b.v=f(b)(v)\qquad\text{et}\qquad b.v=g(b)(v).
\]

On voit que $\mathrm B_i\,\mathrm V_i\subseteq\mathrm V_{i+j}$ et que $\mathrm V_f$ et $\mathrm V_g$ ont même dimension sur $k$ et la proposition 5
nous montre que $\mathrm V_f$ et $\mathrm V_g$ sont des B-Modules isomorphes. On a
\[
s\in\mathrm U(\mathrm A_0)\qquad\text{avec}\qquad s(f(b)(v))=g(b)s(v)
\]
quels que soient $b\in\mathrm B$ et $v\in\mathrm V$ d'où
\[
f(b)(v)=s^{-1}\circ g(b)\circ s(v).
\]

\mspage{11}
